%% kapa: exam=1 round=20 year=2009 period=1 form=A subject=economics subject_ko=경제원론 seq=37 number=77
%% source: Q-Net 감정평가사 기출문제 (cst00309) artlSeq=5024178 file=extracted/1008/5024178/2009년 제20회 감정평가사 1차시험 기출문제/1교시.hwp
%% conversion: kordoc 4.24.51 -> tools/gen1.py (auto); review=proofread-sonnet; flags=none
\begin{question}{77}
순수공공재에 관한 설명으로 옳은 것은?
\begin{choices}
  \item 경합성을 갖는다.
  \item 배제불가능성을 갖는다.
  \item 무임승차자의 문제가 없다.
  \item 정부가 생산하는 모든 재화와 서비스는 순수공공재이다.
  \item 배제불가능성과 경합성을 동시에 갖는 깨끗한 공기는 대표적인 순수공공재이다.
\end{choices}
\end{question}
\answer{2}
